\section*{NS93: averaged cost and the signed correlation that remains}
\paragraph{Scope.} This is an attempted averaged estimate, not a proof of
persistent gain. The descent criterion is prior work. We separate a bounded
diagonal from the actual signed off-diagonal cost and prove that taking
absolute values of all arithmetic products cannot yield the proposed bounded
average. This does not exclude the actual signed sum, the smaller projected
cost, or PR42's full observation route.

Put $N=2^j$, $w_n=\log(2N/n)/n$ and
$T_N(x)=-\sum_{x<n\le2N}\mu(n)w_n$. NS92 gives
$K_N\le U_N\le\kappa\alpha S_N$, where $\alpha=\log2$,
$\kappa=\log(2\pi)-\gamma$, and $S_N=\int_N^{2N}T_N(x)^2dx$.
Finite expansion gives the exact identity
\[
 S_N=D_N+C_N,\quad D_N=\sum_{N<n\le2N}\mu(n)^2w_n^2(n-N),
\]
\[
 C_N=2\sum_{N<m<n\le2N}\mu(m)\mu(n)w_mw_n(m-N).
\]
Since $(n-N)/n^2\le1/(4N)$ and there are $N$ indices,
\[
 0\le D_N\le \alpha^2/4.
\]
Thus an upper bound $\sum_{j=J}^{2J-1} C_{2^j}\le C_*J$ would
supply the cost half of a bounded-average argument. This signed upper bound
is unproved. The diagonal bound supplies no upper bound for $C_N$.

\paragraph{A precise sufficient averaged input, still open.}
Let $b_N$ be NS91's complete one-sided lower numerator, and put
$r_j=(b_{2^j}/E_{2^j})_+$. Where $S_{2^j}>0$, descent gives
\[
 \gamma_j:=\frac{E_{2^j}-E_{2^{j+1}}}{E_{2^j}}
 \ge\frac{r_j^2 E_{2^j}}{\kappa\alpha S_{2^j}}.
\]
A zero direction contributes zero. NS83 supplies $E_{2^j}\ge c/(\alpha j)$
eventually, with $c>0$ and no effective onset here. Consequently
\[
 \sum_{j=J}^{2J-1}\gamma_j
 \ge\frac{c}{2\kappa\alpha^2J}
       \frac{(\sum_{j=J}^{2J-1}r_j)^2}
            {\sum_{j=J}^{2J-1}S_{2^j}}.
\]
This uses weighted Cauchy--Schwarz, retaining the varying actual residuals.
The two hypotheses $\sum r_j\ge\varepsilon J$ and
$\sum S_{2^j}\le C J$, for fixed positive constants on all sufficiently
large dyadic $J$, would therefore give a positive constant per growing
index block and force $E_N\to0$. Neither hypothesis is proved. This is a
less rigid sufficient input than pointwise bounds, not a new convergence
criterion. The index range corresponds to sizes $N$ through $N^2$; it is
not a fixed number of doublings. NS83 still restricts achievable constants.

\paragraph{Absolute products cannot prove that cost hypothesis.}
Define the fully separated majorant
\[
 A_N=\sum_{N<m,n\le2N}|\mu(m)\mu(n)|w_mw_n(\min(m,n)-N).
\]
It is the square integral of the unsigned suffix. For every dyadic
$N\ge2048$,
\[
 A_N\ge\frac{\log^2(4/3)}{9216}N.
\]
Indeed, the interval $(5N/4,3N/2]$ contains exactly $N/4$ integers.
Every nonsquarefree integer is divisible by $k^2$ for some
$2\le k\le\sqrt{3N/2}$. The union bound and
$\sum_{k=2}^\infty k^{-2}\le1/4+\int_2^\infty x^{-2}dx=3/4$
leave at least $N/16-\sqrt{3N/2}\ge N/32$ squarefree integers.
On them $w_n\ge2\log(4/3)/(3N)$; the common overlap length is
at least $N/4$. Taking only these pairs proves the displayed bound.
In particular the absolute off-diagonal allowance $A_N-D_N$ is at least
$N\log^2(4/3)/9216-\alpha^2/4$.

Across a growing dyadic-index block, $\sum_{j=J}^{2J-1}A_{2^j}/J$
therefore diverges. Even the last summand supplies a lower bound proportional
to $2^{2J-1}/J$. Averaging does not repair this particular absolute-product
estimate. No claim of divergence of the signed $S_N$ follows. This is a
failure of a bound, not of the selected direction or of RH.

\paragraph{Finite verification.} Finite suffix and diagonal sums are replayed
at 256/384 bits. A separate direct quadratic expansion checks small sizes;
a doubled sieve limit checks independence from the Mobius table boundary.
There is no physical or infinite tail in these finite identities. Values
and block averages are finite certificates only, not evidence of a uniform
estimate. No optimized residual, new Gram solve, or new lower numerator is
computed in this task.
